\section{Discussion}
\label{sec:discussion}

\paragraph{\textcolor{blue}{Channel and image together}}
\textcolor{blue}{Channel security alone does not decide update security, but the channel still matters (Figure~\ref{fig:channelimage}).
Apple combines TLS~1.3 with signed images whose installation is authorized per device, the model that SUIT and TUF formalize~\cite{rfc9124,samuel2010tuf}.
TP-Link delivers signed images over HTTP, which exposes model and version but not the image's integrity, while its device API, the likely carrier of update commands, lacks forward secrecy.
Reolink's local push of an image without evident protection, and the checksum-only D-Link image, are the cases in which a network adversary's success would depend on the channel alone.
Most fixes are configuration and signing changes rather than hardware changes: enabling forward-secret key exchange, bounding the lifetime of vendor-issued certificates and checking their revocation, and signing images and update metadata as RFC~9124 prescribes.}

\paragraph{\textcolor{blue}{Measurement choices change the conclusions}}
\textcolor{blue}{Two common shortcuts would have reversed several of our findings.
Keyword matching would have analyzed 6,537 captures, 95\% of which contain no update traffic, and it finds some update servers only by accident (Figure~\ref{fig:keyword}).
Statistics over whole captures would have reported long certificate lifetimes that belong to CAs (up to 50~years), weak certificates that belong to the cameras' local servers (self-signed RSA-1024), weak negotiated suites that belong to local app traffic, and an ``encrypted'' verdict for Reolink's unencrypted image.
Separating update traffic first, by the servers involved, is therefore a prerequisite rather than a refinement.}

\paragraph{\textcolor{blue}{Then and now}}
\textcolor{blue}{The 2019 traces show how the same kinds of devices behaved several years earlier, but the devices in them have since been updated, so we compare practices rather than devices.
Between 2019 and our captures, TLS~1.3 appeared on update channels (Apple), and public-CA certificate lifetimes halved, both following ecosystem-wide changes.
Two practices persisted: update checks and images over HTTP (Samsung, LG, and D-Link in 2019; TP-Link today), and long-lived certificates from vendors' own CAs (SmartThings in 2019; TP-Link today).
Only Apple and Amazon appear in both datasets, and their update channels kept authenticated encryption with forward secrecy throughout.}

\paragraph{Limitations and future work}
\textcolor{blue}{(i)~\emph{Scale}: the controlled set has ten devices, of which seven downloaded an update and three were already up to date, and the 2019 traces contain 15 device labels with update traffic from experiments that were not designed to trigger updates; our results characterize practices, not prevalence.
(ii)~\emph{Identification}: we find update servers through name patterns, large transfers, and manual confirmation, so update commands inside multi-purpose device APIs, which we mark as likely, cannot be separated from other traffic within TLS, and an update served by a server whose name and traffic reveal nothing would be missed.
(iii)~\emph{Passive view}: TLS~1.3 hides certificates; we see neither manifests inside TLS nor whether devices check signatures, validate certificates, or pin their vendor's CA; revocation checks outside a capture would be missed; and the image analysis covers only the images we could obtain.
(iv)~\emph{Exposure, not exploitation}: Table~\ref{tab:preconditions} states what each attack requires; we did not attack the devices.
Future work includes active tests of certificate validation and signature checking on the tested devices, and repeated measurements on more devices with scripted update triggers.}
